\chapter{翻译：四种字母怎样变成二十种氨基酸}

\begin{kaichang}
打开你家的冰箱，找三样东西：一块豆腐、一个鸡蛋、一盒牛奶。它们一个来自大豆、一个来自母鸡、一个来自奶牛，颜色、气味、口感天差地别。但如果把它们分别送进实验室做“彻底水解”——把里面的蛋白质完全拆开——你会得到三份几乎一模一样的清单：同样的大约二十种氨基酸，一样不多，一样不少。

反过来也成立。你中午吃了豆腐，大豆蛋白在你的肠道里被拆成氨基酸；几个小时后，这些氨基酸可能已经被你的肌肉细胞重新组装，变成了人的肌肉蛋白。大豆的蛋白质就这样变成了你的蛋白质——原料通用，图纸各归各。

现在猜一猜。细胞用来书写“蛋白质图纸”的信使 RNA 只有四种“字母”，而要指定的氨基酸有二十种。如果每种氨基酸用 1 个字母指定，够用吗？用 2 个字母呢？生命最后选了几个字母一组？多出来的组合是浪费还是另有用处？这一章结束时，我们带着答案回到这块豆腐。
\end{kaichang}

\section{现象：二十种砖块，千万种建筑}

先只做现象层该做的事：描述能测到的事实，不急着解释。

第一个事实：蛋白质的种类多得吓人，原料却少得惊人。你头发里的角蛋白、蛋清里的卵清蛋白、帮你消化食物的酶、对抗病毒的抗体（第 61 章），性质完全不同；但把它们水解到底，得到的都是同一份二十种氨基酸的清单（第 48 章）。而且这份清单不分物种：从大肠杆菌到大象，从蘑菇到你，用的都是同一批砖块。千差万别的是砖块的\textbf{排列顺序}。

第二个事实：蛋白质的生产速度快得离谱。一个大肠杆菌在条件好时二十分钟就分裂一次，这意味着它要在二十分钟内把自己全部的几百万个蛋白质分子重造一份。你的身体也不慢：每天都有几百克旧蛋白质被拆掉、新蛋白质被造出来，肌肉、血浆蛋白、各种酶都在不停更新。

第三个事实有点意外：这台生产机器可以拆下来研究。科学家把细胞轻轻打碎、离心，得到的“无细胞提取物”在试管里仍然能合成蛋白质——只要供给氨基酸和能量。蛋白质合成不是神秘的整体魔法，而是一台可以拆开、可以逐个零件检查的机器。这个试管里的现象，正是后面证据故事的入口。

第四个事实藏在你的餐盘里。二十种氨基酸中，有八九种你的身体自己不会造，必须从食物里现成地吃进来，它们叫“必需氨基酸”。这就是营养学里“蛋白质质量”的全部含义：豆腐厉害的地方不是味道，而是它的氨基酸清单和你的需求对得上号；而长期极端偏食出问题，缺的往往不是“蛋白质”这个总量，而是清单上的某几种砖块——生产线有图纸、有机器，唯独缺料。

第五个事实：序列差一点点，性质就可能差很多。猪和人的胰岛素只差一个氨基酸，所以早年从猪胰腺提取的胰岛素能给人用；而有些遗传病，仅仅是某个蛋白质里换错了一个氨基酸（比如镰刀形贫血症的血红蛋白），整条链的折叠和性质就被改写。蛋白质的“信息”不在原料清单里，而在排列顺序里——这正是“翻译”这个词的分量：顺序必须一个字母一个字母地忠实传达。

\section{粒子层：一条信使、一群搬运工、一座装配台}

把镜头推进到分子。第 68 章讲过，RNA 聚合酶照着 DNA 抄出一份可携带的副本——信使 RNA（mRNA）。这份副本上写着四种核苷酸“字母”：A、U、G、C。可是氨基酸根本不认识碱基。你把一串 mRNA 和一堆氨基酸倒在一起，什么也不会发生：一个说“核苷酸语言”，一个连字母表都没有。两种语言之间，必须有人当翻译。

生命雇了两类角色来解决这个问题：

\begin{itemize}[itemsep=2pt]
\item \textbf{tRNA：一头认字、一头扛货的搬运工。}tRNA 是一小段折成特定形状的 RNA。它的一端露出三个核苷酸，叫\textbf{反密码子}，能按碱基互补规则（第 66 章的 A 对 U、G 对 C）认出 mRNA 上的某三个字母；它的另一端挂着一种特定的氨基酸。它是两种语言之间的活字典。
\item \textbf{氨酰-tRNA 合成酶：给搬运工装货的装卸队长。}氨基酸不会自己爬到正确的 tRNA 上。每一种氨基酸都有专门的一种酶，负责把它装到对应的 tRNA 上，装货要消耗 ATP。这是翻译准确性的第一道关口：装错了货，后面读得再准也白搭。
\end{itemize}

最后登台的是装配台本身：\textbf{核糖体}。它由核糖体 RNA（rRNA）和几十种蛋白质构成，分大小两个亚基，像上下两半模具，把 mRNA 夹在中间。核糖体干的才是翻译的核心工作：它让“拿着氨基酸的 tRNA”依次对号入座，并催化氨基酸之间形成肽键。值得强调一句：核糖体催化的核心部位主要是 RNA，而不是蛋白质——它本质上是一个 RNA 酶。这件事在第 56 章关于生命起源的讨论里是个大线索：会干活的 RNA，可能是比蛋白质更古老的元老。

\begin{center}
\begin{tikzpicture}[scale=1]
\draw[rounded corners, fill=blue!8] (1.6,0.55) rectangle (7.4,2.15);
\node at (4.5,1.95) {\small 核糖体（大亚基）};
\draw[rounded corners, fill=blue!16] (2.2,-0.55) rectangle (6.8,0.45);
\node at (4.5,-0.35) {\small 核糖体（小亚基）};
\draw[thick] (0,0.45) -- (9.6,0.45);
\node[below] at (0.8,0.4) {\small mRNA};
\node at (3.0,0.95) {\small UUU};
\node at (4.5,0.95) {\small AUG};
\node at (6.0,0.95) {\small GGC};
\draw[->, thick] (8.0,0.2) -- (9.3,0.2) node[right] {\small 移动方向};
\draw[fill=green!15] (2.75,1.15) rectangle (3.25,1.55);
\draw[fill=green!15] (4.25,1.15) rectangle (4.75,1.55);
\draw[fill=orange!20] (4.5,1.55) circle (0.16);
\draw[fill=orange!20] (5.15,1.55) circle (0.16);
\draw[thick] (4.66,1.55) -- (4.99,1.55);
\node[right] at (5.35,1.55) {\small 正在加长的肽链};
\node[align=center] at (4.5,-1.3) {\small 示意图：方框、圆圈和颜色只是表示法，不代表真实形状与大小};
\end{tikzpicture}
\end{center}

所以翻译这件事，本质是\textbf{信息的换码}：mRNA 上的核苷酸序列（四种字母），通过 tRNA 这个适配器、在核糖体这台装配台上，被改写成氨基酸序列（二十种砖块）。剩下的问题是：换码的规则是什么？几个字母对应一种氨基酸？

\section{规则：三个字母一个单词}

这是整章最美的一步，因为它只用乘法就能想明白。

假如 1 个字母指定 1 种氨基酸：四种字母最多指定 4 种，不够 20。假如 2 个字母一组：$4\times 4=16$ 种组合，还是不够。3 个字母一组呢？$4\times 4\times 4=64$ 种组合，第一次超过 20。所以\textbf{三联体是能够编码二十种氨基酸的最短方案}。大自然的答案正是 3：mRNA 上连续的三个核苷酸构成一个\textbf{密码子}，指定一种氨基酸。

读法还有三条规则，全是后来实验逼出来的结论（证据故事在后面）：从固定的起点开始读；三个一组、\textbf{不重叠}地连续读；中间没有逗号、没有空格。这意味着起点一旦定错，后面全盘皆输——把阅读框挪动一个字母，就像把“三个和尚没水吃”错读成“个和 尚没 水吃……”，整段话立刻变成乱码。这种现象叫\textbf{移码}，它的破坏力比写错一个字母大得多，第 72 章讲突变时还会回来算账。

\begin{qiaoliang}
64 个格子，只分给 20 种氨基酸加 3 个“停止信号”，多出来的 41 个格子是浪费吗？恰恰相反，这是缓冲。实际情况是：同一种氨基酸常常由好几个密码子指定，比如甘氨酸有 4 个“同义词”，丝氨酸有 6 个——这叫密码的\textbf{简并}。注意简并的规律：同义密码子之间的差别大多在第三个字母上，前两位往往雷打不动。后果很实际：复制或转录时单个字母出错（这类错误总会发生，第 67 章），如果错在第三位，有相当概率根本不改变氨基酸——错误被密码表“吸收”掉了。生命没有追求零错误（那太贵了），而是把代码设计得对错误不敏感。这种“冗余换稳健”的思路，你在工程上会一再见到。
\end{qiaoliang}

密码子里还有两个特殊岗位。一个是\textbf{起始密码子} AUG：它既是指挥部发令枪（从这里开始读、阅读框由它定），又兼任甲硫氨酸的代码——所以每条新造的多肽链，第一个氨基酸几乎都是甲硫氨酸（之后可能被切掉）。另外三个 UAA、UAG、UGA 是\textbf{终止密码子}：它们不指定任何氨基酸，是句号。

\section{符号：密码表登场}

现在可以让符号出场了。完整的遗传密码表有 64 格，任何一本生化教材都能查到；这里只取代表性的一部分，请你观察规律，而不是背格子。

\begin{table}[htbp]
\centering
\begin{tabular}{lll}
\toprule
密码子 & 含义 & 说明 \\
\midrule
AUG & 甲硫氨酸 & 兼任起始信号，定阅读框 \\
UUU、UUC & 苯丙氨酸 & 差别只在第三位 \\
UAU、UAC & 酪氨酸 & 差别只在第三位 \\
GGU、GGC、GGA、GGG & 甘氨酸 & 第三位随便换都行 \\
UCU、UCC、UCA、UCG、AGU、AGC & 丝氨酸 & 同义词最多达 6 个 \\
UAA、UAG、UGA & 终止信号 & 不指定任何氨基酸 \\
\bottomrule
\end{tabular}
\end{table}

配对的样子也写给你看。mRNA 上的密码子从 $5'$ 端读向 $3'$ 端，tRNA 上的反密码子方向相反、碱基互补：

\[ \text{mRNA 密码子}\ 5'\text{-AUG-}3' \quad \longleftrightarrow \quad \text{tRNA 反密码子}\ 3'\text{-UAC-}5' \]

而氨基酸之间的连接，用的是第 48 章认识的肽键：一个氨基酸的氨基与前一个氨基酸的羧基之间脱去一分子水形成的 \chem{-CO-NH-} 连接。在核糖体里，这一步的本质是：旧肽链从旧 tRNA 上“过户”到新 tRNA 带来的氨基酸上，肽键一个接一个地接龙，链就长了。请注意，这一章谈的是信息，但信息从不白来——每一步背后都有第 50 章那套能量账，下面我们就算一算。

顺便教你读完整密码表的方法，以后见到就不会发怵：全表按密码子的第一位、第二位、第三位字母分行分列，查一个密码子就像查地图坐标——先定位第一位，再定位第二位，最后在格子里按第三位找到对应的氨基酸。拿 GGU 练手：第一位 G、第二位 G 锁定“甘氨酸”那一格，第三位无论怎么换都还在格子里，这就是表格排布把“简并”直接画在了脸上。

\section{流水线：开工、延伸、收工}

有了零件和规则，整条流水线可以这样跑：

\textbf{开工（起始）。}核糖体小亚基先搭上 mRNA，滑到第一个 AUG；携带甲硫氨酸的起始 tRNA 靠反密码子对号入座；大亚基从上面扣下来，装配台合拢。从这一刻起，阅读框被锁定。

\textbf{延伸（循环）。}接下来是三步循环，每转一圈加一个氨基酸：① 进位——新 tRNA 带着氨基酸进来试座，反密码子配不上就被弹出去（这是第二道校对）；② 成键——核糖体催化肽键形成，旧肽链整个“过户”到新来的氨基酸上；③ 移位——核糖体沿 mRNA 挪动三个字母，空出货位，腾出下一个座位。

\textbf{收工（终止）。}当终止密码子滑进座位，没有任何 tRNA 认领它——细胞本来就没给这三个密码子配搬运工。来报到的是一种叫\textbf{释放因子}的蛋白质，它让整条多肽链脱离核糖体，大小亚基随之解体，零件回收，等待下一条 mRNA。

这套流水线准吗？相当准，但绝不是零误差。装货酶抓错氨基酸、反密码子认错密码子的事都会发生，层层校对之后，翻译的总错误率大约是每掺入一万个氨基酸出错一次。拿它和第 67 章的 DNA 复制比一比：复制的错误率低到十亿分之一量级。为什么两道工序的精度差这么多？想想代价就明白了：一条多肽造错了，报废的只是这一条链，细胞再造一条就是；DNA 抄错了，错误会留在图纸原本上，传给这个细胞的全体后代。校对是要花能量、花时间的，生命在每道工序上的精度投入，恰好和“出错的代价”相称——这又是能量与概率那条老规则在管事。

这条流水线快吗？细菌的核糖体每秒大约能加 20 个氨基酸，你的细胞里每秒也有好几个。一条 300 个氨基酸的多肽，一个核糖体十几秒到一分钟就能造完；而且一条 mRNA 上常常串着好几个核糖体同时施工，像一根铁丝上穿着一串珠子，这叫多聚核糖体。一个细胞里成千上万台装配台连轴转，才撑得起前面说的“每天更新几百克蛋白质”。

\begin{qiaoliang}
算一笔诚实的能量账。往肽链上加一个氨基酸，大约要花掉 4 个“ATP 当量”：装货时 ATP 被拆到 AMP（相当于 2 个），进位和移位各用 1 个 GTP。一个成年人每天合成的蛋白质按约 $300$ 克估算，氨基酸残基平均约 $110$ 道尔顿，折合每天形成约 $2.7$ 摩尔肽键，也就是约 $1.6\times 10^{24}$ 个。乘上 4，再按每摩尔 ATP 水解放出约 $50$ 千焦换算，你一天花在蛋白质合成上的能量大约是 $500$ 千焦上下——约占一天总能耗的百分之几。也就是说，哪怕你整天躺着，你的身体也在为“信息变成物质”持续付费。这个估算的每一步你都可以质疑和改写：合成量、能量单价，换上你自己查到的数字再算一遍。
\end{qiaoliang}

\section{证据：密码是怎样被破译的}

\begin{zhengju}
1953 年 DNA 双螺旋公布后，所有人都意识到下一个问题：四种字母，怎么指定二十种氨基酸？

最先认真推算的反而是物理学家伽莫夫。他算出了“至少三联体”这道乘法题，还提出一个漂亮的猜想：密码是重叠的，相邻密码子共用字母。这个猜想有一个可检验的预言——如果密码重叠，一种氨基酸换了，它旁边的氨基酸就会受牵连，蛋白质的氨基酸排列就会出现特定限制。等科学家测出的蛋白质序列多起来，发现根本没有这种限制：猜想被干净利落地否决了。理论再漂亮，也要给数据让路。

与此同时，克里克在 1955 年提出一个更大胆的预言：氨基酸不可能直接认碱基，细胞里一定存在一种“适配器”分子，一头认密码、一头挂氨基酸。当时没有任何人见过这种分子。几年后，tRNA 被发现，形状和功能与预言惊人吻合——和门捷列夫的空格子一样，先预言、后找到，是科学理论最硬气的时刻。

真正打开密码表的是实验。1961 年，克里克和布伦纳用一类能在 DNA 里插入或删除单个碱基的化学药剂做遗传实验：插 1 个字母，基因报废；插 2 个，还是报废；插 3 个——功能竟然大体恢复了！这正是“三个字母一组、从固定点连续读”的直接证据。同年，尼伦伯格和马太把人工合成的、只含 U 的 RNA（poly-U）放进无细胞蛋白质合成体系，结果试管里只造出一种“蛋白质”：一串纯粹的苯丙氨酸。第一个密码子就此破译：UUU = 苯丙氨酸。对照实验很清楚——不加这种 RNA，就什么都不合成。随后几年，科拉纳等人合成出各种已知顺序的重复 RNA（比如 UCUCUC……，造出丝氨酸和亮氨酸交替的链），一格一格地把整张表填满。到 1966 年，64 个密码子全部确定。

请留意这个故事的形状：先有数学推算，再有被否决的猜想、被证实的预言，最后靠设计巧妙的对照实验逐格填空。没有一步靠“天才的直觉”直接跳到答案。
\end{zhengju}

\section{边界：翻译完成，不等于蛋白质完工}

这一章的模型——密码表加翻译流水线——回答的是“氨基酸按什么顺序排”。它回答不了“这条链长成什么形状、干什么活、去哪个岗位”。一条刚从核糖体里放出来的多肽链，多半还不能工作：

\begin{itemize}[itemsep=2pt]
\item \textbf{要折叠。}肽链要按第 48 章讲的规律折成特定的立体形状，有时还需要“分子伴侣”蛋白质帮忙；折错的链会被细胞标记、拆解回收。
\item \textbf{要修饰。}许多蛋白质要被剪掉一段（胰岛素就是先造成一条长链，再切去中间一段才有活性）、接上糖链、打上磷酸基团（第 58 章讲过磷酸化怎样当开关用），才算成品。
\item \textbf{要运输。}肽链前端常常带有一段“信号肽”，像快递单上的地址，决定它被送去线粒体、细胞膜还是分泌到细胞外。地址错了，蛋白质再好也到不了岗位。
\end{itemize}

还有一条信息层面的边界：密码表只管“是什么”，不管“有多少”。同一条 mRNA 可以被翻译一次，也可以被翻译上千次；蛋白质的实际数量取决于 mRNA 有多少、活多久、翻译多快——这些密码表一概不说。这个缺口，正是通往下一章的门。

\begin{shiyan}
\textbf{在纸上当一次核糖体。}

材料：纸、笔，以及本章的密码子片段表。

步骤：① 在纸上抄下这条 mRNA 序列（从 $5'$ 端读）：AUGGGUUUUCACUAA；② 从第一个 AUG 开始，每三个字母画一条分隔线，逐格查表，写出氨基酸顺序；③ 现在做一次“突变”：把第 6 个字母 G 改成 A，重新翻译，比较前后的氨基酸序列；④ 再做一次更狠的：在第 4 个字母后面插入一个 A，重新分组、重新翻译，看看从哪一位开始全乱、乱到哪里。

观察点：第 ③ 步中，GGU 变成 GGA 后氨基酸变没变？这就是简并的缓冲作用。第 ④ 步中，一个多余的字母为什么比改错一个字母破坏力大得多？

安全提示：纯纸笔活动，唯一的危险是数错格子——数错了请从头再数，核糖体可不会数错。
\end{shiyan}

\begin{wuqu}
“遗传密码是一张全宇宙通用、永远不变、绝无例外的密码表。”——前半句几乎对：从细菌到大象确实共用同一本字典，这本身就是所有生命同出一源的有力证据。但“绝无例外”是错的。你的线粒体（那座远古细菌的后代）就私自改了几条规则：在细胞质里 UGA 是终止信号，在人的线粒体里它却编码色氨酸。更妙的是，UGA 在特定场合还能编码第二十一种氨基酸——硒半胱氨酸，你的抗氧化酶里就有它；某些纤毛虫干脆把终止密码子改派了别的用场。密码表不是物理定律，它是演化的产物，而演化的产物总会留下改动的痕迹。把这个误解当回事的不只是考试：做基因工程时，如果把基因按“标准表”直接搬进密码规则不同的系统里表达，产物就会出错——第 84 章还会遇到这类坑。
\end{wuqu}

\section{回到开场的豆腐}

现在可以兑现承诺了。开场问题的答案是：1 个字母只够 4 种，2 个字母只够 16 种，都不够 20；3 个字母给出 64 格，第一次够用——所以生命用三联体密码。多出来的格子不是浪费：3 格当了句号，其余变成“同义词”，为拼写错误提供缓冲。

豆腐、鸡蛋、牛奶为什么那么不同？因为大豆、母鸡、奶牛的 mRNA 序列不同，密码子读出的氨基酸排列就不同；排列不同，折叠出的形状就不同；形状不同，口感、气味、营养就不同。而它们为什么又能互相“通用”——你吃了豆腐能长出人的肌肉？因为三家的氨基酸原料是同一批二十种，密码字典也是同一本。你的肠道只认砖块不认建筑：拆掉别人的图纸，用自己的 mRNA 重新指挥核糖体施工。你和一块豆腐之间，隔着的是信息，不是物质。

\section{装配台之外}

这套机制一清楚，几件大事立刻有了着落。

先说医学。你药箱里的红霉素软膏、医生开的四环素，杀的正是核糖体：这些抗生素专门卡住细菌核糖体的某个部位，让它的翻译流水线停摆。你的核糖体结构和细菌的差异足够大，同样剂量下基本不受影响——这就是“杀菌不伤身”的分子基础。当然，剂量和个体差异永远重要，而且抗生素分不清坏菌和你肠道里的好房客（第 62 章），滥用还会筛选出抗药的细菌——那是自然选择送上门来的账单（第 76 章）。

反过来看，核糖体也是你的命门。白喉杆菌之所以致命，就是因为它释放的毒素能给你的核糖体“上锁”，让细胞的翻译全线瘫痪。一条流水线对所有蛋白质都不可替代——这是效率的来源，也是脆弱的来源。

再说演化线索。你的线粒体里有一套自己的核糖体，无论大小还是对抗生素的敏感性，都更像细菌的而不像你细胞质里的——线粒体是远古细菌后代的证据，就刻在它的装配台上（第 56 章的内共生故事）。

还有工程。既然密码表近乎通用，把人胰岛素的基因搬进细菌，细菌就能老老实实地翻译出人胰岛素——今天糖尿病人用的胰岛素大多这样生产，第 83 章会细讲这条路怎样走通。

但留一个更大的问题作伏笔：这一章的流水线随叫随到，可一个细胞有几万个基因，难道全都开足马力翻译到底？神经细胞和肌肉细胞带着同一套 DNA，凭什么造出完全不同的蛋白质阵容？谁决定哪些 mRNA 被造出来、造多少、活多久？装配台只是执行者，指挥部另有其人——第 70 章，我们去看基因是怎样被打开、关小和关掉的。而密码表里的“同义词”，还埋着另一笔账：DNA 改了一个字母，蛋白质有时居然纹丝不动——这种“沉默的改变”到底是错误、材料还是创新，第 72 章见分晓。

\begin{tiaozhan}
简并到底能吸收多少错误？你可以亲手算出来。任取一个密码子，它有 9 种“改一个字母”的改法（3 个位置，每个位置有 3 种改法）。以 GGU（甘氨酸）为例：改第三位（GGU$\rightarrow$GGC、GGA、GGG）全都不变；改第一位或第二位，几乎必然换成别的氨基酸。把 64 个密码子的 $64\times 9=576$ 种单字母改变全部查一遍表，你会发现大约四分之一完全不改变氨基酸——这还没算“换成性质相似的氨基酸”这类温和改变。换句话说，密码表在第三位上装了减震器。再进一步：为什么减震恰好装在第三位？提示——想想 tRNA 反密码子与密码子配对时，哪一位的碱基配对要求最松（这叫“摆动”）。跳过本节不影响主线。
\end{tiaozhan}

\section*{练一练}

\begin{sikaoti}
\item 画一张信息流图：从 DNA 到成品蛋白质，标出 mRNA、tRNA、氨酰-tRNA 合成酶、核糖体各自的位置，并在每个箭头旁注明“这一步说的是哪种语言、谁负责换码、能量花在哪一步”。
\item 用本章的密码子片段表翻译序列 AUGUAUGGCUAA（从 $5'$ 端读），写出氨基酸顺序；然后把阅读框向后挪一位重新分组，描述发生了什么。
\item 一个细菌核糖体每秒约加 20 个氨基酸。一条 500 个氨基酸的多肽，单个核糖体需要多久？按每加一个氨基酸消耗 4 个 ATP 当量，造一条这样的多肽大约要多少 ATP 当量？
\item 假如某种外星生命用 5 种核苷酸字母、2 个字母一组编码 20 种氨基酸，理论上够用吗？与我们“4 种字母、3 个一组”的方案相比，各有什么得失？（提示：从组合总数、密码子总数与简并缓冲三个角度比较。）
\item 红霉素卡住细菌核糖体却不直接卡死你的细胞，依据是什么？但长期服用广谱抗生素的人常会肠胃不适——用第 62 章的生态系统观点解释为什么。
\item 设计一个实验，验证“AAA 编码赖氨酸”这一格密码表：你会用什么“人工 mRNA”、什么反应体系、观察什么产物、设什么对照？
\end{sikaoti}
